% !TeX root = ../main.tex
\begin{abstract}
\noindent\textbf{Background.}  Beta-binomial models represent a guide count as
a proportion conditional on its library total.  Existing CRISPR implementations
of this idea are largely restricted to two-group comparisons, whereas modern
screens often require time, dose, donor, or interaction effects.

\noindent\textbf{Methods.}  BARCS implements guide-level beta-binomial
regression,
$\logit(\mu_{gi})=\mathbf{x}_i^\top\boldsymbol{\beta}_g$.  It conditions each
guide count on the unfiltered library total, estimates prespecified
coefficients or contrasts, and optionally uses negative-control calibration
and dispersion moderation.  Simulations with known truth and held-out controls
were used to distinguish ranking performance from error-rate calibration.

\noindent\textbf{Results.}  In a correctly specified continuous-dose
simulation, unmoderated BARCS was anti-conservative: its null type-I error was
0.081 and its empirical false-discovery proportion was 0.130 at nominal 0.05,
versus 0.038 and 0.049 for MAGeCK-MLE.  A five-fold held-out analysis of 593
non-targeting guides in an IL2RA screen gave a BARCS $p<0.05$ rate of 0.049
after control scaling, but the correlated FACS-bin design remains outside the
independent-library model.  A Cas13 sensitivity analysis used days 0, 7, and 14
in four replicate-complete cell lines; BARCS had a lower mean
cell-line-specific null calibration
error than MAGeCK-MLE (0.025 versus 0.057), whereas MAGeCK-MLE and general
RNA-sequencing methods ranked essential genes more strongly.  Because these
inputs were normalized, ComBat-corrected, and rounded, that comparison does not
validate either count likelihood.  Dispersion and denominator ablations
improved BARCS in simulations, but they do not establish superiority over
competing methods.

\noindent\textbf{Conclusions.}  BARCS supplies a design-matrix extension of
beta-binomial CRISPR-screen analysis, but the present plug-in dispersion
estimator does not provide uniformly calibrated finite-sample inference.
Independent biological replication, likelihood-compatible counts, and
held-out or cross-fitted controls are required.  The results support the
software contribution and define the work needed for confirmatory use; they do
not show that beta-binomial regression dominates negative-binomial methods.
\end{abstract}
